\section{Formal reduction method and scope of guarantees}
\label{sec:formal-method}
This formulation retains the label-distribution weights, directed dissimilarity and deterministic granular-ball construction. It specifies prediction feasibility as the reduction constraint and granular dependency as search guidance. Acquisition cost refers to measuring selected attributes at prediction time; it is distinct from algorithmic running time.

\subsection{Native LDL representation and granular guidance}
Let $T=\{(x_i,\boldsymbol\ell_i)\}_{i=1}^n$, where $\boldsymbol\ell_i\in[0,1]^q$ and $\sum_{j=1}^q\ell_{ij}=1$. Let $A$ contain $m$ attributes with finite costs $c_a>0$, and define $C(P)=\sum_{a\in P}c_a$. No single dominant decision label replaces the native distribution.

The entropy-based information weight $W_a(i)$ measures the divergence between a regularized attribute-local average distribution and the global training distribution. Its interpretation is local supervised association, rather than guaranteed predictive importance. The directed dissimilarity combines this information with raw attribute geometry:
\begin{equation}
 \rho_P(i,j)=\frac{1}{m}\left[(1-\zeta)\sum_{a\in P}W_a(i)\Delta_a(i,j)
 +\zeta\sum_{a\in P}G_a(i,j)\right],\quad 0<\zeta<1.
\end{equation}
Here $\Delta_a$ is the normalized numeric discrepancy or the total variation between categorical label profiles; $G_a$ is the numeric discrepancy or categorical mismatch. All component statistics are fixed when comparing subsets within a fitting sample. The denominator is the original $m$, not $|P|$. Thus $P\subseteq Q$ implies $\rho_P(i,j)\le\rho_Q(i,j)$. Source-dependent weights allow asymmetry; no triangle inequality is required.

Granular construction uses $d_P(i,j)=\max\{\rho_P(i,j),\rho_P(j,i)\}$, keeps identical selected-attribute vectors together, and applies deterministic farthest-seed splitting, support checks and fixed eligibility rules. Let $\mathcal E_P$ denote the eligible leaves, $\boldsymbol p=n^{-1}\sum_i\boldsymbol\ell_i$, and $V_0=n^{-1}\sum_i\|\boldsymbol\ell_i-\boldsymbol p\|_2^2$. For $V_0>0$,
\begin{equation}
 \Gamma(P)=\frac{\sum_{B\in\mathcal E_P}|B|\|\boldsymbol\mu_B-\boldsymbol p\|_2^2}{nV_0}
 =1-\frac{R_{\mathrm{GB,tr}}(P)}{V_0},\qquad 0\le\Gamma(P)\le1.
\end{equation}
$R_{\mathrm{GB,tr}}$ is the mean squared training representation error: members of eligible balls use their ball mean, while all remaining members use $\boldsymbol p$. Set $\Gamma(P)=0$ when $V_0=0$. The identity follows from the within--between decomposition on each eligible leaf. It concerns training representation, not predictions on unseen objects. Regenerating balls after adding attributes may decrease $\Gamma$; signed gains are therefore retained.

\subsection{Prediction-feasible minimum-cost objective}
Partition the development training sample into fixed inner validation folds $I_1,\ldots,I_K$. For each subset, refit preprocessing and the predictor using the complementary fitting rows. Define
\begin{equation}
 R_{\mathrm{inner}}(P)=\frac{1}{n}\sum_{k=1}^K\sum_{i\in I_k}
 \|\boldsymbol\ell_i-\widehat{\boldsymbol\ell}^{(-k)}_P(x_i)\|_2^2,
 \qquad L_{\max}=R_{\mathrm{inner}}(A).
\end{equation}
The reference predictor is the mean label distribution of five Euclidean nearest fitting neighbours after fitting-sample min--max scaling. The empty subset predicts the fitting-sample mean distribution. A predictor using supervised geometry would additionally require all label-derived weights and profiles to be fitted separately within every fold.

The reduction objective is
\begin{equation}
 \mathcal F=\{P\subseteq A:R_{\mathrm{inner}}(P)\le L_{\max}\},
 \qquad P^*\in\operatorname*{argmin}_{P\in\mathcal F}C(P).
 \label{eq:formal-cost-objective}
\end{equation}
The full set is feasible by construction. If the empty set is feasible, it is the unique zero-cost optimum. The historical constraint $\Gamma(P)\ge\Gamma(A)$ is a comparator, not part of this prediction-feasibility formulation: training representation and predictor loss need not impose compatible orderings.

\subsection{Dependency-guided constructive search}
For $a\notin P$, write
\begin{equation}
 u_a(P)=R_{\mathrm{inner}}(P)-R_{\mathrm{inner}}(P\cup\{a\}),\qquad
 v_a(P)=\Gamma(P\cup\{a\})-\Gamma(P).
\end{equation}
Use the established cost exponent $\lambda=-1$, so positive gains are ranked per unit cost. The following steps define the base heuristic; ties use smaller acquisition cost and then smaller attribute index.
\begin{enumerate}
\item Initialize $P=\varnothing$. If $P\in\mathcal F$, return $P$.
\item While $P\notin\mathcal F$, evaluate $u_a(P)$ for all remaining attributes. Let $J=\{a:u_a(P)>0\}$ and $H=\{a\in J:v_a(P)>0\}$.
\item If $H\ne\varnothing$, choose the maximum $v_a(P)/c_a$ over $H$. Otherwise, if $J\ne\varnothing$, choose the maximum $u_a(P)/c_a$ over $J$. If $J=\varnothing$, choose the maximum raw signed $u_a(P)$ over all remaining attributes. Add the selected attribute and repeat the feasibility check.
\item Scan selected attributes in decreasing cost order, with increasing index for ties. Accept any deletion preserving membership in $\mathcal F$ and restart the scan. Stop after a complete unsuccessful scan.
\end{enumerate}
The nonpositive bridge permits continued construction when single additions do not improve the current loss. It does not guarantee recovery of every useful interaction. Dependency ranks a restricted set of prediction-improving candidates without introducing a mixing coefficient. This priority can select a smaller predictive improvement, so its benefit is empirical and may depend on the budget and dataset.

\begin{proposition}[Base heuristic guarantee]
Assuming finite deterministic subset losses, the base heuristic terminates with a feasible, single-deletion-irreducible subset after at most $m$ additions and $m$ accepted deletions. It need not return an inclusion-minimal or globally minimum-cost subset.
\end{proposition}
\begin{proof}
Each addition introduces an unused attribute. If earlier subsets fail, the construction reaches the feasible full set. Each accepted deletion reduces cardinality and preserves feasibility. A final complete unsuccessful scan excludes every feasible single deletion. Nonmonotone loss allows a multiple deletion to be feasible even when each single deletion fails, so stronger minimality does not follow.
\end{proof}
There are $O(m^2)$ subset evaluation requests in the base search. This bound counts requests, not the computational cost of fitting predictors or constructing balls.

\subsection{Optional lower-cost exchange refinement}
For a feasible subset $P$, consider deletions and exchanges $Q=(P\setminus\{a\})\cup\{b\}$ with $a\in P$, $b\notin P$ and $c_b<c_a$. Order candidates by decreasing cost saving, removed index and added index. Evaluate candidates subject to a predeclared request budget. Among evaluated feasible candidates, select the one with smallest cost, then smallest inner loss, then the move indices. Accept it and restart. The diagnostic implementation uses a total budget of 200 requests, including cache hits. A partial scan can accept its best evaluated feasible move before stopping at the budget.

Every accepted move strictly reduces acquisition cost and preserves the full-set loss guard. With finitely many subsets, unrestricted refinement terminates. Only termination after a complete unsuccessful neighbourhood scan establishes local optimality against these deletion and cheaper single-exchange moves. Budget exhaustion provides no such guarantee and may leave a feasible deletion after an accepted exchange. Loss may increase relative to the preceding subset while staying below $L_{\max}$. This phase is an explicit optional algorithm revision, not a consequence of the base heuristic theorem.

\subsection{Exact reference and limits of the claims}
An include/exclude enumeration with incumbent $A$ solves Eq.~\eqref{eq:formal-cost-objective} exactly when completed. Prune a node whose current cost is at least the incumbent cost. A feasible current subset can replace the incumbent, and its strict extensions need not be searched because all added costs are positive. No dependency- or loss-monotonicity pruning is permitted. Every omitted extension is therefore no cheaper than an available feasible solution; finite enumeration yields a global minimum. The worst-case search is exponential. No completed exact optimum on the high-dimensional real datasets is claimed.

The exact objective is independent of $\Gamma$: granular information changes the heuristic path, not the feasible family. The contributions supported by this formulation are a native LDL information geometry, a distribution-based granular representation, and reproducible cost-sensitive search with an internal prediction constraint despite nonmonotone dependency. Claims of uniformly superior prediction, lossless external-test reduction, or universally beneficial granular guidance are not established.

Inner folds are reused adaptively during search. Membership in $\mathcal F$ is an internal selection guarantee and does not imply an unbiased validation estimate or external predictive preservation. Independent held-out predictions and common-budget comparisons assess empirical cost--loss tradeoffs. Previously inspected test splits remain exploratory evidence. The directed geometry is used for selection and granulation in the reference pipeline; the Euclidean predictor is stated explicitly rather than silently replaced by a directed-neighbour predictor.
